\chapter{Dataset, analisi esplorativa e preprocessing}\label{ch:dati}
\section{Struttura e attività degli utenti}
Last.fm 1K raccoglie eventi acquisiti attraverso \emph{user.getRecentTracks}; il README locale dichiara una riga per ascolto e sei campi separati da tab: utente, timestamp, artist MBID, nome artista, track MBID e titolo \cite{lastfm}. Il file completo contiene 19.150.868 righe e 992 utenti. Il profilo anagrafico separato non entra nelle feature del modello. Gli eventi sono osservazioni implicite di ascolto, non giudizi espliciti né click su raccomandazioni del sistema.

Il notebook \metric{01_dataset_analysis.ipynb} distingue un campione iniziale di 100.000 righe dalla scansione esplorativa completa. Il campione contiene soltanto quattro utenti: le sue percentuali non sono quindi trattate come stime rappresentative dell'intera popolazione. Le statistiche complete mostrano una forte eterogeneità: gli eventi per utente hanno media 19.305,31, mediana 11.551, minimo 2 e massimo 183.103. Il primo e il terzo quartile sono 3.774,25 e 25.707,75. Si sceglie la macro-media per assegnare pari peso agli utenti; una media per decisione attribuirebbe più peso ai profili con più sessioni e rappresenterebbe un diverso obiettivo di aggregazione.

\figura{eda_activity.png}{.91}{Distribuzione degli ascolti per utente: numero di eventi sull'asse orizzontale e utenti in scala logaritmica su quello verticale. Output originale della cella 21 del notebook EDA, ingrandito per leggibilità; nessuna nuova scansione.}{fig:activity}

Le entità esplorative contate nel notebook sono 176.892 artisti e 1.505.209 tracce. Tali numeri dipendono dai criteri di identificazione usati nell'EDA e non rappresentano automaticamente il numero di item canonici della pipeline definitiva. In particolare, l'identità operativa finale include l'artista anche quando esiste un track ID. Analogamente, il catalogo per la coverage è la sola unione dei candidati esposti nella fase: usare il conteggio globale EDA come denominatore cambierebbe la metrica.

\clearpage
\section{Concentrazione, ripetizione e sparsità}
La popolarità è molto concentrata. Nel conteggio esplorativo, l'1\% degli artisti raccoglie il 62,7254\% degli ascolti e il 10\% il 91,9006\%. Per le tracce le quote corrispondenti sono 31,8500\% e 71,1172\%. La mediana della popolarità è di quattro eventi per artista e due per traccia. Questi dati descrivono una lunga coda con molti contenuti poco supportati. Non è disponibile in questa relazione una misura di densità della matrice canonica utente--item; la sparsità è motivata dall'eterogeneità e dalla coda osservate, senza inventare una percentuale di celle vuote.

\figura{eda_concentration.pdf}{.91}{Concentrazione degli ascolti nell'EDA. Curve ridisegnate dai cinque valori tabulati già salvati nelle celle 39--40 del notebook: nessun ricalcolo dal dataset. La curva degli artisti corrisponde alla Figura originale della cella 41.}{fig:concentration}

Anche la ripetizione individuale è marcata: il tasso medio esplorativo di ripetizione è 0,885690 per artista e 0,664859 per traccia; le mediane sono 0,944116 e 0,724130. Il 10\% degli artisti preferiti da ciascun utente concentra mediamente il 62,9290\% dei suoi ascolti. Questa evidenza giustifica una baseline personalizzata forte e la distinzione fra nuova traccia di artista noto e nuovo artista.

\clearpage
\section{Qualità dei dati e identità canoniche}
L'EDA rileva 602.166 eventi senza artist ID (3,144327\%) e 2.168.588 senza track ID (11,323706\%). Scartarli significherebbe rimuovere una frazione sostanziale della storia. Il loader mantiene i campi vuoti e costruisce fallback espliciti. Il nome esatto dell'artista sostituisce il suo ID soltanto quando questo manca; il titolo sostituisce il track ID soltanto se assente. Maiuscole, spazi significativi e Unicode non sono normalizzati arbitrariamente; un valore composto solo da spazi è considerato mancante. Se non è possibile identificare artista o item, viene sollevato un errore.

Il file è TSV, non CSV con quoting: ogni riga fisica possiede cinque tab. Il parser definitivo usa \metric{QUOTE_NONE}, così le virgolette nei titoli restano letterali. Il controllo a sei colonne rimane attivo. Il precedente comportamento CSV poteva fondere più righe, spiegando perché il notebook preliminare \metric{reward_sessions_split.ipynb} riporti 19.098.853 righe e coorti di 561 utenti. Quegli output servono come motivazione esplorativa, non sostituiscono i 19.150.868 eventi e i 562 utenti delle run ufficiali.

Un secondo problema riguarda track ID attribuiti a più artisti. In record di MF DOOM e Viktor Vaughn compare lo stesso ID:
\begin{quote}\small\path{7aebbc02-17ab-4f94-b91d-ff164c89f1f8}.\end{quote}
L'identità finale è uniforme per tutto il dataset: se presente, il track ID è abbinato all'artist key; altrimenti il titolo esatto è abbinato alla stessa artist key. Tag di tipo e lunghezze dei campi rendono la serializzazione non ambigua. La regola non dipende da collisioni scoperte a posteriori o dalla frequenza futura dell'item.

\section{Tempo, sessioni e motivazione delle soglie}
L'EDA mostra storie molto diverse per durata: la mediana è circa 925,63 giorni. Sono presenti anche timestamp isolati nel 2010 e nel 2013; non vengono corretti inventando date, ma i limiti delle fasi impediscono che entrino nelle metriche sperimentali del 2009.

Il notebook preliminare sui gap osserva che, fra pause inferiori a due ore, il 94,89\% ricade entro dieci minuti. Nei riascolti della stessa traccia, la finestra di sette giorni raccoglie circa il 50,98\% dei gap misurati, contro il 29,98\% a un giorno e il 60,54\% a quattordici. Sono indicazioni esplorative ottenute prima della correzione del parser: motivano le soglie adottate, ma non sono stime aggiornate della reward artist-level.

Nel codice definitivo una sessione inizia al primo evento dell'utente oppure quando il gap dall'evento precedente è \emph{strettamente maggiore} di dieci minuti. Il decision point coincide con quell'evento, che resta escluso dallo stato usato dalla policy. Il caso gap esattamente pari a dieci minuti non apre una nuova sessione. Gli split e l'eleggibilità finale sono formalizzati nel Capitolo~\ref{ch:metodo}.

\paragraph{Implicazioni progettuali dell'analisi esplorativa.}
Concentrazione e ripetizione rendono necessario confrontarsi con una baseline personalizzata forte, ma suggeriscono anche di misurare quante proposte escano dalla familiarità storica. L'eterogeneità motiva il pari peso agli utenti, mentre gli ID mancanti richiedono identità utilizzabili senza perdere eventi. La sequenza degli ascolti impone di aggiornare storia e candidati nel tempo. Le analisi preliminari dei gap orientano sessionizzazione e orizzonte, senza dimostrarne l'ottimalità. Queste osservazioni motivano i requisiti dell'architettura seguente; non costituiscono un confronto sperimentale fra architetture alternative.
